\section{Why the channels differ}
\label{sec:why-channels-differ}

Three measurements are consistent with, though they do not prove, a read-vs-not-read
account of the dissociation. A linear probe trained to decode the active belief from
hidden states succeeds when compiled modulation is active and degrades to 22\% accuracy
when it is zeroed, indicating the belief's functional signature is concentrated in the
modulated pathway rather than diffused. Attention-pattern divergence between modulated
and unmodulated forward passes on the same prompt is small (Jensen--Shannon divergence
0.049), consistent with modulation acting on the value/query content of attention rather
than substantially redirecting where the model attends. And modulation measurably
suppresses think-mode generation, the pathway through which the model would otherwise
narrate --- and thereby expose to an interlocutor's rebuttal --- its own reasoning about
whether to concede.

We state the claim at the strength the evidence supports: the two channels are
\emph{dissociable}, and the compiled channel affects processing without appearing in the
reasoning trace the model produces. This is consistent with, not proof of, an account in
which context-delivered content is read and is therefore available to be argued against,
while compiled content is not. We flag directly that the same probe result above is used
in Appendix~\ref{app:probe} to make a different point: decodability from hidden states,
even at reduced accuracy, is enough to fail a strict implicit-memory criterion. Both
readings use the same number and do not contradict each other --- the belief is
decodable by a probe with access to internals, and it is nonetheless never surfaced in
the text the model reasons over or the interlocutor sees.
